% mode: ini
% Aimed at after_math's section 1379, each way through: the display and its number too wide (w+q>z),
% with shrink enough to keep the number, with too little, so the number goes on a line of its own, with
% infinite shrink of each order; the formula alone too wide, squeezed by hpack; and the boundary, in
% scaled points, of w+q>z and w-shrink+q<=z, with q=e+\quad of \textfont2 at 10pt.  Also section 1381's
% l and e=0, and 1383's: the number on a line by itself, before and after.  The overfull messages.
\input prelude
\def\e#1{\vrule width#1pt height1pt}
\def\sh{\showbox1 }
\def\d#1{\setbox1\vbox{\noindent$$#1$$}\sh}
\hbadness=0 \hfuzz=0pt \overfullrule=0pt
\hsize=100pt
\abovedisplayskip=11pt \belowdisplayskip=12pt \abovedisplayshortskip=13pt \belowdisplayshortskip=14pt
% q is e+18pt or so, whichever \quad is: show it, then use it.
\dimendef\qd=200
\setbox1\hbox{\qd=\fontdimen6\tensy \global\qd=\qd}
\showthe\qd
% Room, exactly, and 1sp too little: the number goes down without shrink.
\dimen0=100pt \advance\dimen0 by-\qd \advance\dimen0 by-20pt
\d{\vrule width\dimen0 height 1pt\eqno\e{20}}
\advance\dimen0 by 1sp
\d{\vrule width\dimen0 height 1pt\eqno\e{20}}
\d{\vrule width\dimen0 height 1pt\leqno\e{20}}
% With finite shrink: w-shrink+q=z exactly, and 1sp more.
\dimen1=1sp
\d{\vrule width\dimen0 height 1pt\hskip 0pt minus\dimen1\eqno\e{20}}
\d{\vrule width\dimen0 height 1pt\hskip 0pt minus 0.5\dimen1\eqno\e{20}}
\d{\vrule width\dimen0 height 1pt\kern 5pt\hskip 0pt minus 5pt\eqno\e{20}}
\d{\vrule width\dimen0 height 1pt\kern 5pt\hskip 0pt minus 4.99999pt\eqno\e{20}}
\d{\vrule width\dimen0 height 1pt\kern 5pt\hskip 0pt minus 4pt\eqno\e{20}}
\d{\vrule width\dimen0 height 1pt\kern 5pt\hskip 0pt minus 4pt\leqno\e{20}}
\d{\vrule width\dimen0 height 1pt\kern 5pt\hskip 0pt minus 6pt\eqno\e{20}}
\d{\vrule width\dimen0 height 1pt\kern 5pt\hskip 0pt minus 6pt\leqno\e{20}}
% Infinite shrink: always kept, however wide.
\d{\e{300}\hskip 0pt minus 1fil\eqno\e{20}}
\d{\e{300}\hskip 0pt minus 1fill\eqno\e{20}}
\d{\e{300}\hskip 0pt minus 1filll\eqno\e{20}}
\d{\e{300}\hskip 0pt minus 1fil\leqno\e{20}}
\d{\e{300}\hskip 0pt plus 1fil\eqno\e{20}}
% Number wider than the line.
\d{\e{10}\eqno\e{200}}
\d{\e{10}\leqno\e{200}}
\d{\eqno\e{200}}
\d{\e{10}\hskip 0pt minus 1fil\eqno\e{200}}
% No number, the formula too wide: with and without shrink; z and w are compared without q.
\d{\e{100}}
\d{\vrule width 100.00002pt height 1pt}
\d{\e{300}}
\d{\e{110}\hskip 0pt minus 5pt}
\d{\e{110}\hskip 0pt minus 15pt}
\d{\e{110}\hskip 0pt minus 1fil}
% Width 0: e is not 0, w+q>z by the number's width alone.
\displaywidth=0pt
\setbox1\vbox{\noindent$$\displaywidth=5pt \eqno\e{3}$$}\sh
\setbox1\vbox{\noindent$$\displaywidth=0pt \e{5}$$}\sh
\setbox1\vbox{\noindent$$\displaywidth=-5pt \e{5}$$}\sh
\setbox1\vbox{\noindent$$\displaywidth=-5pt \e{5}\eqno\e{3}$$}\sh
\setbox1\vbox{\noindent$$\displaywidth=-5pt \eqno$$}\sh
% Hanging, and \displayindent, moved.
\setbox1\vbox{\noindent$$\displayindent=9pt \e{300}\eqno\e{20}$$}\sh
\setbox1\vbox{\noindent$$\displayindent=9pt \e{300}\leqno\e{20}$$}\sh
\setbox1\vbox{\noindent$$\displayindent=9pt \e{10}\leqno\e{20}$$}\sh
% Show the skips of a number on its own line, with the code g2 going to 0.
\belowdisplayskip=0pt \belowdisplayshortskip=0pt
\d{\e{300}\eqno\e{20}}
\d{\e{10}}
\belowdisplayskip=12pt \belowdisplayshortskip=14pt
\end
